%%%%%%%%%%%%%%%%%
% CV tailored for Schneider Electric - Consultant Analytique Energie
% Positioning: Consultant Analytique Energie / Data Scientist Industriel | ML, IoT, performance, optimisation & reporting
%%%%%%%%%%%%%%%%%

\documentclass[8pt,a4paper,ragged2e,withhyper]{altacv}

\geometry{left=1.25cm,right=1.25cm,top=.5cm,bottom=1.cm,columnsep=1.2cm}

\usepackage{paracol}
\usepackage{fontawesome5}
\usepackage{hyperref}

\iftutex
  \setmainfont{Roboto Slab}
  \setsansfont{Lato}
  \renewcommand{\familydefault}{\sfdefault}
\else
  \usepackage[rm]{roboto}
  \usepackage[defaultsans]{lato}
  \renewcommand{\familydefault}{\sfdefault}
\fi

\definecolor{SlateGrey}{HTML}{2E2E2E}
\definecolor{LightGrey}{HTML}{666666}
\definecolor{DarkPastelRed}{HTML}{8AC0D9}
\definecolor{PastelRed}{HTML}{00B0DA}
\definecolor{GoldenEarth}{HTML}{B4AD0C}

\colorlet{name}{black}
\colorlet{tagline}{PastelRed}
\colorlet{heading}{DarkPastelRed}
\colorlet{headingrule}{GoldenEarth}
\colorlet{subheading}{PastelRed}
\colorlet{accent}{PastelRed}
\colorlet{emphasis}{SlateGrey}
\colorlet{body}{LightGrey}

\definecolor{violet}{HTML}{5A2582}
\definecolor{rouge}{HTML}{F53B3B}
\definecolor{noir}{HTML}{00232C}
\definecolor{bleu}{HTML}{00B0DA}
\definecolor{rose}{HTML}{F831A0}
\definecolor{jaune}{HTML}{FFEC00}
\definecolor{vert}{HTML}{34AD6C}

\renewcommand{\namefont}{\Huge\rmfamily\bfseries}
\renewcommand{\personalinfofont}{\footnotesize}
\renewcommand{\cvsectionfont}{\LARGE\rmfamily\bfseries}
\renewcommand{\cvsubsectionfont}{\large\bfseries}
\renewcommand{\cvItemMarker}{{\small\textbullet}}
\renewcommand{\cvRatingMarker}{\faCircle}

\input{../../_shared/pubs-num.tex}

% auto_tex_manager layout tuning
\emergencystretch=3em
\hfuzz=60pt
\hbadness=9999
\sloppy

\begin{document}

\name{BOILEAU Guillaume}
\tagline{Consultant Analytique Energie / Data Scientist Industriel | ML, IoT, performance, optimisation \& reporting}

\photoL{2.8cm}{../../_shared/moi}

\personalinfo{%
  \email{guillaume.boileaupro@gmail.com}
  \phone{+33 6 59 94 85 84}
  \location{Alpes-Maritimes, FRANCE}
  \linkedin{guillaume-boileau-891b81265}
  \researchgate{Guillaume-Boileau-2}
  \orcid{0000-0002-3576-6968}
}

\makecvheader
\vspace{-0.3cm}

\AtBeginEnvironment{itemize}{\small}
\columnratio{0.64}

\begin{paracol}{2}

\cvsection{Experience}

\cvevent{\textbf{Software Engineer -- Données opérationnelles, mobilité électrique \& support terrain}}{VEV Platform Services France}{2025 -- 2026}{Valbonne, France}

\textbf{Contexte :} Plateforme CPMS connectée à des infrastructures de recharge de véhicules électriques, avec services applicatifs, flux de données opérationnels, contraintes terrain, exploitation et enjeux de fiabilité.

\textbf{Objectif :} Contribuer à la fiabilisation de services numériques liés à l'énergie et à la mobilité électrique, en analysant les flux, les incidents, les comportements applicatifs et les contraintes opérationnelles.

\begin{itemize}
  \item \textbf{Énergie \& mobilité électrique :} Travail dans l'écosystème des infrastructures de recharge, à l'interface entre logiciel, données opérationnelles et contraintes terrain.
  \item \textbf{Analyse de données opérationnelles :} Investigation de flux FTP, interruptions d'échanges, incohérences de données et comportements inattendus.
  \item \textbf{Diagnostic \& amélioration continue :} Analyse d'incidents, identification de causes, formalisation de constats et contribution aux actions correctives.
  \item \textbf{Validation opérationnelle :} Vérification de la cohérence entre données applicatives, comportement des équipements et attendus fonctionnels.
  \item \textbf{Développement logiciel :} Contribution à des composants TypeScript, Angular et Node.js intégrés dans une plateforme opérationnelle.
  \item \textbf{Reporting technique :} Rédaction de constats, analyses d'anomalies et supports de résolution pour les équipes.
  \item \textbf{Outils collaboratifs :} Utilisation de Git, Linux, Zenhub et documentation technique pour suivre les sujets.
\end{itemize}

\divider

\cvevent{\textbf{Ingénieur R\&D senior -- Data science, ML, performance \& validation}}{CNES}{2023 -- 2025}{Nice, France}

\textbf{Contexte :} Développement logiciel et R\&D pour le segment sol de la mission spatiale LISA ESA/NASA, avec chaînes de données mission L0--L3, simulation, intégration de modèles, validation, performance et environnement CNES/ESA.

\textbf{Objectif :} Développer des workflows Python robustes pour intégrer, comparer, valider et exploiter des modèles et données complexes, avec exigences de reproductibilité, qualité de données, performance et documentation.

\begin{itemize}
  \item \textbf{Python \& data pipelines :} Développement de CATpyLISA, plateforme Python pour structurer, comparer, valider et documenter des modèles.
  \textcolor{DarkPastelRed}{\href{https://gitlab.oca.eu/gboileau/lisa_l3_tools}{\bf GitLab: CATpyLISA}}
  \item \textbf{Performance \& indicateurs :} Structuration de métriques, contrôles de cohérence, comparaisons de configurations et analyse des écarts.
  \item \textbf{Validation de modèles :} Définition de critères d'acceptation, règles de comparaison, procédures de vérification et analyse d'anomalies.
  \item \textbf{Simulation \& faisabilité :} Développement d'environnements numériques permettant d'évaluer la robustesse des modèles selon différents scénarios.
  \item \textbf{Qualité des données :} Traçabilité des entrées/sorties, versions, hypothèses, limites méthodologiques et résultats.
  \item \textbf{Reporting \& recommandations :} Production de synthèses, supports de revue, documentation technique et éléments d'aide à la décision.
  \item \textbf{Coordination :} Interface avec contributeurs scientifiques, logiciels et institutionnels dans un environnement international.
\end{itemize}

\divider

\cvevent{\textbf{Data Scientist -- Données capteurs, ML \& optimisation de performance}}{Universiteit Antwerpen, Belgium}{2021 -- 2023}{Antwerp, Belgium}

\textbf{Contexte :} Travaux de data science pour instruments de précision, combinant données capteurs, bruit environnemental, modélisation statistique, machine learning, validation de performance et recommandations techniques.

\textbf{Objectif :} Développer des workflows robustes pour analyser des données capteurs, identifier les facteurs limitant la performance, comparer des configurations et proposer des améliorations.

\begin{itemize}
  \item \textbf{Données capteurs / IoT-like :} Analyse de mesures environnementales, séries temporelles, propagation, corrélations spatiales et impacts sur la performance système.
  \item \textbf{Machine learning appliqué :} Approches de réduction de bruits non stationnaires avec canaux témoins, machine learning et deep learning.
  \item \textbf{Pipelines Python :} Développement de workflows pour analyse de données, statistiques, validation, visualisation et reporting.
  \item \textbf{Optimisation de performance :} Figures de mérite, analyses de sensibilité, contrôles de cohérence et comparaison de configurations.
  \item \textbf{Modélisation statistique :} Pipelines MCMC pour différentes configurations instrumentales et différents niveaux de bruit.
  \item \textbf{Recommandations :} Production de synthèses techniques et éléments de décision pour parties prenantes internationales.
\end{itemize}

\divider

\cvevent{\textbf{Ingénieur R\&D junior -- Simulation, signal \& analyse quantitative}}{ARTEMIS, Observatoire de la Côte d'Azur}{2018 -- 2021}{Nice, France}

\textbf{Contexte :} Recherche appliquée liée à l'analyse de signaux faibles, à la simulation numérique, au traitement du signal et à la préparation de missions spatiales.

\textbf{Objectif :} Développer des méthodes d'analyse et de simulation afin de produire des résultats robustes, documentés et exploitables.

\begin{itemize}
  \item \textbf{Traitement du signal :} Méthodes pour analyser et séparer des signaux faibles et superposés.
  \item \textbf{Simulation numérique :} Développement d'approches de simulation pour environnements complexes et grands jeux de données.
  \item \textbf{Analyse quantitative :} Workflows Python/C d'analyse, visualisation, comparaison et validation de résultats.
  \item \textbf{Validation :} Vérification de cohérence, comparaison de résultats, études de sensibilité et analyse de limites méthodologiques.
  \item \textbf{Rédaction :} Formalisation des méthodes, hypothèses, résultats et conclusions pour publications et revues collaboratives.
\end{itemize}

\switchcolumn

\cvsection{Profile}

\textbf{Data Scientist / Consultant Analytique Energie avec expérience en ML, données capteurs, performance, validation de modèles, reporting et environnement énergie/mobilité électrique.}

Positionnement pour ce rôle : analyser des données IoT historiques et temps réel, évaluer la faisabilité d'optimisations énergétiques par ML, gérer les cas d'usage de la définition au déploiement, suivre les performances des modèles et formuler des recommandations exploitables pour les équipes énergie, projet et client.

\begin{itemize}
  \item Solide expérience en Python, ML, analyse de données, séries temporelles, données capteurs, validation et performance.
  \item Expérience concrète en environnement énergie/mobilité électrique via VEV et infrastructures de recharge.
  \item Capacité à transformer des données complexes en diagnostics, KPI, recommandations et supports de décision.
  \item Culture industrielle : services opérationnels, données terrain, incidents, contraintes d'exploitation et amélioration continue.
  \item Montée ciblée sur plateformes ML low-code/no-code, AVEVA/Schneider, systèmes utilities et quantification d'économies d'énergie.
\end{itemize}

\cvsection{Languages}

\cvskill{English}{4}
\divider
\cvskill{Spanish}{2}
\divider

\medskip

\cvsection{Education}

\cvevent{PhD in Physics}{Université Côte d'Azur}{2018 -- 2021}{Nice, France}

\divider

\cvevent{Selective Graduate Program in Fundamental Physics (Magistère)}{Université Paris-Saclay}{2016 -- 2018}{Orsay, France}

\divider

\cvevent{MSc in Astronomy and Astrophysics}{Observatoire de Paris / Université Paris-Saclay}{2017 -- 2018}{Paris, France}

\divider

\cvevent{BSc in Fundamental Physics}{Université Paris-Saclay}{2015 -- 2016}{Orsay, France}

\cvsection{Professional Training}

\cvevent{Machine Learning in Python with scikit-learn}{FUN MOOC -- Inria}{2026}{France Université Numérique}

\small{
Predictive modelling, supervised learning, model evaluation, cross-validation, metrics, pipelines and methodological/software fundamentals of machine learning in Python.
}

\divider

\cvevent{Recherche reproductible : principes méthodologiques pour une science transparente}{FUN MOOC -- Inria}{2025}{France Université Numérique}

\small{
Reproducible research, GitLab, notebooks/Jupyter, data management, computational documentation, software environments and reproducibility methods.
}

\divider

\cvevent{Continuous Professional Training -- Data, AI, Software and Systems}{OpenClassrooms / Coursera / FUN MOOC}{2020 -- 2026}{}

\small{
Python, SQL, Machine Learning, AI, Linux, JavaScript, TypeScript, Angular, Node.js, Django, REST APIs, data analysis, software engineering and algorithms.
}

\end{paracol}

\cvsection{Projects}

\cvevent{VEV -- Mobilité électrique, données opérationnelles \& optimisation de services}{CPMS / EV charging / Operational data}{2025 -- 2026}{}

Contribution à une plateforme numérique connectée à des infrastructures de recharge de véhicules électriques, avec analyse de flux, incidents, données opérationnelles, contraintes terrain et fiabilisation de service.

\textbf{Rôle :} analyse de flux, investigation d'incidents, validation de comportements applicatifs, documentation technique et contribution à l'amélioration continue.

\textbf{Apport Schneider :} compréhension concrète d'un environnement énergie, données terrain, infrastructure connectée, diagnostic opérationnel et qualité de service.

\divider

\cvevent{Noise reduction and predictive performance workflows}{Machine learning / Sensor data / Performance optimisation}{2021 -- 2023}{}

Développement de workflows numériques et statistiques pour analyser des données capteurs, réduire des bruits non stationnaires, comparer des configurations et améliorer la performance de systèmes complexes.

\textbf{Rôle :} analyse de données, ML, modélisation statistique, pipelines MCMC, validation de performance, KPI, recommandations et reporting.

\textbf{Apport Schneider :} logique d'optimisation fondée sur les données, surveillance de performance, validation de modèles et amélioration continue.

\divider

\cvevent{CATpyLISA -- Plateforme Python de modèles, données \& validation}{Mission spatiale LISA ESA/NASA -- CNES/ESA}{2023 -- 2025}{}

Plateforme Python dédiée à l'intégration de modèles, à la structuration de données, à la comparaison de résultats, à la validation technique et à la revue collaborative de workflows scientifiques.

\textbf{Rôle :} développement Python, structuration de données, validation de modèles, analyse d'écarts, documentation technique et coordination avec plusieurs contributeurs.

\textbf{Apport Schneider :} pipelines réutilisables, qualité des données, reproductibilité, validation de modèles et transformation d'analyses en livrables exploitables.

\cvsection{Compétences clés}

\vspace{-.3cm}

\cvsubsection{Data science, ML \& optimisation}

{\small
\cvtag{Machine Learning}
\cvtag{Data Science}
\cvtag{Modèles prédictifs}
\cvtag{Validation de modèles}
\cvtag{Analyse d'erreurs}
\cvtag{Performance monitoring}
\cvtag{Amélioration continue}
\cvtag{Optimisation}
\cvtag{Études de faisabilité}
\cvtag{Benchmarking}
\cvtag{KPI}
\cvtag{Recommandations}
}

\divider\smallskip
\vspace{-.3cm}

\cvsubsection{Énergie, IoT \& industrie}

{\small
\cvtag{Données IoT}
\cvtag{Données capteurs}
\cvtag{Séries temporelles}
\cvtag{Performance énergétique -- ramp-up}
\cvtag{Mobilité électrique}
\cvtag{Infrastructures de recharge}
\cvtag{Systèmes industriels}
\cvtag{Utilities -- ramp-up}
\cvtag{Air comprimé -- awareness}
\cvtag{Pompage -- awareness}
\cvtag{Ventilation / AHU -- awareness}
\cvtag{Vapeur / eau chaude -- awareness}
}

\divider\smallskip
\vspace{-.3cm}

\cvsubsection{Pipelines, plateformes \& low-code}

{\small
\cvtag{Pipelines de données}
\cvtag{Workflows ML}
\cvtag{Déploiement modèles -- ramp-up}
\cvtag{Plateformes IoT -- ramp-up}
\cvtag{Low-code ML -- ramp-up}
\cvtag{No-code ML -- ramp-up}
\cvtag{AVEVA -- ramp-up}
\cvtag{Bibliothèques ML réutilisables}
\cvtag{Data quality}
\cvtag{Reproductibilité}
\cvtag{Monitoring}
}

\divider\smallskip
\vspace{-.3cm}

\cvsubsection{Python, data \& outils}

{\small
\cvtag{Python}
\cvtag{pandas}
\cvtag{NumPy}
\cvtag{SciPy}
\cvtag{scikit-learn}
\cvtag{Matplotlib}
\cvtag{Jupyter}
\cvtag{SQL}
\cvtag{Power BI}
\cvtag{Excel}
\cvtag{Git / GitLab}
\cvtag{Linux}
\cvtag{Docker}
\cvtag{CI/CD}
}

\divider\smallskip
\vspace{-.3cm}

\cvsubsection{Projet, Agile \& collaboration}

{\small
\cvtag{Gestion de cas d'usage ML}
\cvtag{Définition à déploiement}
\cvtag{Suivi planning}
\cvtag{Communication obstacles}
\cvtag{Animation -- ramp-up}
\cvtag{Agile -- pratique}
\cvtag{Scrum / SAFe -- awareness}
\cvtag{Équipes matricielles -- ramp-up}
\cvtag{Coordination transverse}
\cvtag{Clients / parties prenantes}
}

\divider\smallskip
\vspace{-.3cm}

\cvsubsection{Communication \& reporting}

{\small
\cvtag{Analyse \& synthèse}
\cvtag{Rapports techniques}
\cvtag{Présentations client}
\cvtag{Vulgarisation}
\cvtag{Documentation}
\cvtag{Aide à la décision}
\cvtag{Recommandations opérationnelles}
\cvtag{Anglais courant}
\cvtag{Innovation}
\cvtag{Durabilité}
}

\end{document}
